%%% FIELD proof-selected-pair-strict-gap
Write vectors in the coordinate order \(0,1,2\).  By \(\mathrm{def:finite-witness}\) and the exact marginalization in \(\mathrm{thm:finite-witness-nonattainment}\),
\[
B_0=\begin{pmatrix}
109/800&57/800&3/100\\
57/800&9/200&27/800\\
3/100&27/800&39/800
\end{pmatrix},\qquad
B_1=\begin{pmatrix}
39/800&27/800&3/100\\
27/800&9/200&57/800\\
3/100&57/800&109/800
\end{pmatrix},
\]
\[
b_0=\begin{pmatrix}21/400\\9/200\\21/400\end{pmatrix},\qquad
b_1=\begin{pmatrix}3/50\\21/200\\37/200\end{pmatrix},\qquad
d=\begin{pmatrix}7/20\\3/10\\7/20\end{pmatrix}.
\]
The displayed selected representatives are
\[
h_0^{\mathrm{sel}}=\begin{pmatrix}0\\2/5\\4/5\end{pmatrix},\quad
h_1^{\mathrm{sel}}=\begin{pmatrix}1/5\\3/5\\1\end{pmatrix},\quad
q_0^{\mathrm{sel}}=\begin{pmatrix}0\\8/3\\16/3\end{pmatrix},\quad
q_1^{\mathrm{sel}}=\begin{pmatrix}16/3\\8/3\\0\end{pmatrix}.
\]
Direct multiplication gives the four identities
\[
B_0h_0^{\mathrm{sel}}=b_0,\qquad
B_1h_1^{\mathrm{sel}}=b_1,
\qquad
B_0^\top q_0^{\mathrm{sel}}=d,\qquad
B_1^\top q_1^{\mathrm{sel}}=d.
\]
Indeed, the first two products have numerators, after use of a common denominator \(4000\), equal respectively to
\[
(210,180,210)^\top/4000=b_0,qquad
(240,420,740)^\top/4000=b_1,
\]
and the last two products both equal \((7/20,3/10,7/20)^\top\).

We now translate these unconditional moment equations into the asserted conditional bridge equations.  Put
\(r_a(z)=P^\star(A=a,Z=z)=\sum_wB_a[z,w]\).  Positivity of all \(36\) observed cells, proved in \(\mathrm{thm:finite-witness-nonattainment}\), implies \(r_a(z)>0\) and \(d(w)=P^\star(W=w)>0\).  By the definitions of \(T_a\), \(g_a\), \(B_a\), and \(b_a\),
\[
(T_ah)(z)=\frac{(B_ah)[z]}{r_a(z)},
\qquad
g_a(z)=\frac{b_a[z]}{r_a(z)}.
\]
Consequently \(B_ah_a^{\mathrm{sel}}=b_a\) implies
\(T_ah_a^{\mathrm{sel}}=g_a\).  Similarly, the definition of the adjoint gives
\[
(T_a^*q)(w)
=E_{P^\star}[I_aq(Z)\mid W=w]
=\frac{(B_a^\top q)[w]}{d(w)}.
\]
Thus \(B_a^\top q_a^{\mathrm{sel}}=d\), together with \(d(w)>0\), implies
\(T_a^*q_a^{\mathrm{sel}}=1\).  This proves all four bridge restrictions without invoking uniqueness.

For the variance calculation, \(h_1^{\mathrm{sel}}(w)-h_0^{\mathrm{sel}}(w)=1/5\) for every \(w\), and \(\mathrm{thm:finite-witness-nonattainment}\) gives \(\theta(P^\star)=1/5\).  Hence the centered selected score reduces to
\[
\phi_{\mathrm{sel}}
=I_1q_1^{\mathrm{sel}}(Z)\{Y-h_1^{\mathrm{sel}}(W)\}
-I_0q_0^{\mathrm{sel}}(Z)\{Y-h_0^{\mathrm{sel}}(W)\}
=t_Aq_A^{\mathrm{sel}}(Z)\{Y-h_A^{\mathrm{sel}}(W)\}.
\]
Under \(\mathrm{def:finite-witness}\), conditional on \((A=a,U^\star=u)\), the variables \(Z\), \(W\), and \(Y\) are mutually independent, the law of each of \(Z,W\) is
\[
p_0=(3/5,3/10,1/10),\qquad p_1=(1/10,3/10,3/5),
\]
and
\[
\mu_{a,u}:=P^\star(Y=1\mid A=a,U^\star=u)=\frac{1+a+2u}{5}.
\]
Therefore
\[
E[(q_a^{\mathrm{sel}}(Z))^2(Y-h_a^{\mathrm{sel}}(W))^2\mid A=a,U^\star=u]
=Q_{a,u}H_{a,u},
\]
where
\[
Q_{a,u}=\sum_{z=0}^2p_u(z)(q_a^{\mathrm{sel}}(z))^2,
\qquad
H_{a,u}=\sum_{w=0}^2p_u(w)
\left[\mu_{a,u}\{1-h_a^{\mathrm{sel}}(w)\}^2
+(1-\mu_{a,u})\{h_a^{\mathrm{sel}}(w)\}^2\right].
\]
Substitution of the displayed vectors gives the following exact table:
\[
\begin{array}{c|c|c|c|c|c}
u&a&P^\star(U^\star=u,A=a)&Q_{a,u}&H_{a,u}&P^\star(U^\star=u,A=a)Q_{a,u}H_{a,u}\\ \hline
0&0&3/8&224/45&29/125&812/1875\\
0&1&1/8&96/5&39/125&468/625\\
1&0&1/8&96/5&39/125&468/625\\
1&1&3/8&224/45&29/125&812/1875
\end{array}
\]
Since \(I_0I_1=0\), summing the last column proves
\[
V_{\mathrm{sel}}
=E_{P^\star}[\phi_{\mathrm{sel}}^2]
=2\frac{812}{1875}+2\frac{468}{625}
=\frac{4432}{1875}.
\]
The established witness theorem gives
\(V_{\mathrm{eff}}(P^\star)=773723312/334026875\).  Exact subtraction yields
\[
V_{\mathrm{sel}}-V_{\mathrm{eff}}(P^\star)
=\frac{4432}{1875}-\frac{773723312}{334026875}
=\frac{9496288}{200416125}>0,
\]
which proves the claimed strict comparison at \(P^\star\).

It remains to justify persistence, including the asserted gauge construction.  Work in the relatively open positive rank-two joint-solvability chart supplied by \(\mathrm{def:finite-constraint-map}\) and \(\mathrm{thm:finite-witness-nonattainment}\).  At \(P^\star\), both \(\ker(B_a)\) and \(\ker(B_a^\top)\) are spanned by
\(v=(3,-7,3)^\top\), so every coordinate used below is nonzero in the relevant null vector.

For the outcome bridge in arm \(a\), impose the gauge
\[
h_a(P)(0)=h_a^{\mathrm{sel}}(0)=a/5.
\]
Because \(v(0)\ne0\) and \(B_a(P^\star)\) has rank two, the two columns indexed by \(w=1,2\) are linearly independent.  Hence there is a pair of row indices \(\mathscr I_a\) for which
\(B_a(P^\star)[\mathscr I_a,\{1,2\}]\) is nonsingular.  On the neighborhood where this determinant remains nonzero, define
\[
\begin{pmatrix}h_a(P)(1)\\h_a(P)(2)\end{pmatrix}
=B_a(P)[\mathscr I_a,\{1,2\}]^{-1}
\left{b_a(P)[\mathscr I_a]
-B_a(P)[\mathscr I_a,\{0\}]\frac a5\right}.
\]
This is a rational, hence smooth, function of the finite cell probabilities.  Outcome solvability and rank two imply that the remaining row equation also holds.  The gauge intersects each affine solution fiber once because every nonzero vector in \(\ker(B_a(P))\) has a nonzero zeroth coordinate after the neighborhood is shrunk.

For treatment bridges let \(r_0=0\) and \(r_1=2\), and impose
\[
q_a(P)(r_a)=q_a^{\mathrm{sel}}(r_a)=0.
\]
Since \(v(r_a)\ne0\), the two rows of \(B_a(P^\star)\) with indices different from \(r_a\) are linearly independent.  Choose two columns on which the resulting two-by-two minor is nonsingular and solve the corresponding two equations of
\(B_a(P)^\top q_a(P)=d(P)\) for the two unrestricted coordinates.  The resulting formula again consists only of matrix inversion and multiplication by finite moment coordinates, so it is rational and smooth while that minor remains nonzero.  Treatment solvability and rank two give the remaining equation.  These four sections agree with the displayed selected representatives at \(P^\star\) by uniqueness under their gauges.

For completeness, choose rationally varying null vectors \(u_a(P)\in\ker(B_a(P))\) and \(m_a(P)\in\ker(B_a(P)^\top)\) by fixing one coordinate that equals \(3\) at \(P^\star\).  Form the six normal generators
\[
D(P)=\bigl(F_0(h_0(P),m_0(P)),G_0(q_0(P),u_0(P)),R_0(m_0(P),u_0(P)),
F_1(h_1(P),m_1(P)),G_1(q_1(P),u_1(P)),R_1(m_1(P),u_1(P))\bigr)^\top.
\]
Because both kernels are one-dimensional, these six functions span \(\mathcal C_P\).  Let
\[
\Gamma(P)=E_P[D(P)D(P)^\top],qquad
c(P)=E_P[D(P)\phi_{\mathrm{sel}}(P)].
\]
The exact base-law Gram calculation in \(\mathrm{thm:finite-witness-nonattainment}\) is
\[
\Gamma(P^\star)=
\begin{pmatrix}
1323/500&-6/5&63/20&0&0&0\\
-6/5&343/5&-21&0&-21&0\\
63/20&-21&441/2&0&0&0\\
0&0&0&1323/500&6/5&-63/20\\
0&-21&0&6/5&343/5&-21\\
0&0&0&-63/20&-21&441/2
\end{pmatrix},
\qquad
\det\Gamma(P^\star)=\frac{16253767998768573}{12500000}>0.
\]
As a Gram matrix it is nonnegative definite, and its positive determinant makes it positive definite.  All entries of \(D(P)\), \(\Gamma(P)\), \(c(P)\), and the selected-score variance are continuous rational functions of the positive cell probabilities on the chosen chart.  Thus, after shrinking the chart, \(\Gamma(P)\) remains invertible.

The affine-normal identity used by \(\mathrm{thm:finite-witness-nonattainment}\) gives
\(\mathcal A_P=\phi_{\mathrm{sel}}(P)+\mathcal C_P\).  Hence finite-dimensional orthogonal projection onto the span of \(D(P)\) gives
\[
V_{\mathrm{eff}}(P)
=E_P[\phi_{\mathrm{sel}}(P)^2]-c(P)^\top\Gamma(P)^{-1}c(P),
\]
and therefore the continued selected-pair gap is
\[
\Delta(P):=E_P[\phi_{\mathrm{sel}}(P)^2]-V_{\mathrm{eff}}(P)
=c(P)^\top\Gamma(P)^{-1}c(P).
\]
This is continuous, and at the base law
\(\Delta(P^\star)=9496288/200416125>0\).  A final shrinking of the relatively open positive rank-two chart therefore ensures \(\Delta(P)>0\) throughout it.  This proves the stated smooth rational continuation and persistence of the strict variance gap.  The argument compares only the displayed bridge-score section with the quotient efficiency bound and makes no assertion about the probability limit of any published estimator.
